\chapter{Peaks over threshold}

The Pickands--Balkema--de Haan theorem (Balkema and de Haan, 1974; Pickands, 1975) says
that, for distributions in a maximum domain of attraction, the conditional distribution of
the excess over a high threshold $u$ is, after rescaling, close to a generalized Pareto
distribution. We prove the heavy-tailed (Fréchet) case; see Embrechts, Klüppelberg and
Mikosch, \emph{Modelling Extremal Events} (1997), Theorem~3.4.13 [EKM], and McNeil, Frey
and Embrechts, \emph{Quantitative Risk Management} (2005), \S5.2 [MFE].

Throughout, $F \colon \bR \to \bR$ plays the role of a survival function
$F(x) = \PR[X > x]$.

\begin{definition}[Generalized Pareto survival function]
  \label{def:gpd-survival}
  \lean{ExtremeValueTheory.gpdSurvival}
  \leanok
  For shape $\xi \ge 0$, scale $\beta > 0$ and $x \ge 0$, the survival function of the
  generalized Pareto distribution (GPD) is
  \begin{align*}
    \overline{G}_{\xi,\beta}(x) \; = \;
    \begin{cases}
      (1 + \xi x / \beta)^{-1/\xi} & \text{if } \xi \ne 0, \\
      \exp(-x/\beta) & \text{if } \xi = 0 .
    \end{cases}
  \end{align*}
\end{definition}

\begin{lemma}[Basic properties of the GPD]
  \label{lem:gpd-properties}
  \uses{def:gpd-survival}
  \lean{ExtremeValueTheory.gpdSurvival_of_eq_zero, ExtremeValueTheory.gpdSurvival_of_ne_zero,
    ExtremeValueTheory.gpdSurvival_pos, ExtremeValueTheory.gpdSurvival_le_one,
    ExtremeValueTheory.gpdSurvival_antitoneOn, ExtremeValueTheory.gpdSurvival_tendsto_atTop,
    ExtremeValueTheory.gpdSurvival_continuousOn, ExtremeValueTheory.gpdSurvival_inv_one}
  \leanok
  Let $\xi \ge 0$ and $\beta > 0$. On $[0, \infty)$ the function $\overline{G}_{\xi,\beta}$ takes
  values in $(0, 1]$, is antitone and continuous, and tends to $0$ at $+\infty$. Moreover, for
  $\alpha \ne 0$, $\overline{G}_{1/\alpha, 1}(x) = (1 + x/\alpha)^{-\alpha}$.
\end{lemma}
\begin{proof}
  \leanok
  For $x \ge 0$ the base $1 + \xi x / \beta$ is at least $1$, and the exponent $-1/\xi$ is negative,
  so all properties follow from elementary properties of real powers and of $\exp$.
\end{proof}

\begin{definition}[Excess survival function and mean excess function]
  \label{def:excess-survival}
  \lean{ExtremeValueTheory.excessSurvival, ExtremeValueTheory.meanExcess}
  \leanok
  The \emph{conditional excess survival function} over the threshold $u$ is
  $\overline{F}_u(x) = F(u + x)/F(u)$; if $F(x) = \PR[X > x]$, this is $\PR[X - u > x \mid X > u]$.
  The \emph{mean excess function} is $e(u) = \big( \int_u^\infty F(t) \, \ud t \big) / F(u)$, which
  equals $\EX[X - u \mid X > u]$ when $\EX[X^+] < \infty$.
\end{definition}

\begin{lemma}[Pointwise convergence of excesses]
  \label{lem:pot-pointwise}
  \uses{def:regularly-varying, def:gpd-survival, def:excess-survival, lem:gpd-properties}
  \lean{ExtremeValueTheory.IsRegularlyVarying.tendsto_excessSurvival}
  \leanok
  If $F$ is regularly varying with index $-\alpha$, $\alpha > 0$, then with
  $\beta(u) = u/\alpha$, for each $x \ge 0$,
  \begin{align*}
    \lim_{u \to \infty} \frac{F(u + \beta(u) x)}{F(u)} \; = \; \overline{G}_{1/\alpha, 1}(x) .
  \end{align*}
\end{lemma}
\begin{proof}
  \leanok
  We have $u + \beta(u) x = (1 + x/\alpha) u$, so by regular variation with $c = 1 + x/\alpha > 0$
  the ratio tends to $(1 + x/\alpha)^{-\alpha} = \overline{G}_{1/\alpha,1}(x)$
  (Lemma~\ref{lem:gpd-properties}).
\end{proof}

\begin{theorem}[Pickands--Balkema--de Haan, Fréchet case]
  \label{thm:pickands-balkema-de-haan}
  \uses{lem:pot-pointwise, thm:rv-uct-monotone, lem:gpd-properties}
  \lean{ExtremeValueTheory.IsRegularlyVarying.tendstoUniformlyOn_excessSurvival,
    ExtremeValueTheory.IsRegularlyVarying.tendsto_sup_excessSurvival}
  \leanok
  Let $F$ be antitone and regularly varying with index $-\alpha$, $\alpha > 0$, and let
  $\beta(u) = u/\alpha$. Then $x \mapsto F(u + \beta(u) x)/F(u)$ converges to
  $\overline{G}_{1/\alpha, 1}$ uniformly on $[0, \infty)$ as $u \to \infty$; that is,
  \begin{align*}
    \lim_{u \to \infty} \; \sup_{x \ge 0} \; \Big| \frac{F(u + \beta(u) x)}{F(u)}
      - \overline{G}_{1/\alpha, 1}(x) \Big| \; = \; 0
  \end{align*}
  [EKM, Theorem~3.4.13].
\end{theorem}
\begin{proof}
  \leanok
  Let $\eps > 0$. Since $\overline{G}_{1/\alpha,1}(x) \to 0$, choose $M$ with
  $\overline{G}_{1/\alpha,1}(M) < \eps/2$. On the compact interval $c = 1 + x/\alpha \in [1, 1 + M/\alpha]$
  the uniform convergence theorem for monotone regularly varying functions
  (Theorem~\ref{thm:rv-uct-monotone}) gives uniform convergence for $x \in [0, M]$. For $x > M$,
  both functions are nonnegative and bounded by their values at $x = M$ (by monotonicity), which
  are eventually less than $\eps$; hence the difference is less than $\eps$ there as well.
  The supremum form is a reformulation of uniform convergence.
\end{proof}

\begin{theorem}[Mean excess function of a regularly varying tail]
  \label{thm:mean-excess}
  \uses{def:excess-survival, thm:karamata}
  \lean{ExtremeValueTheory.IsRegularlyVarying.tendsto_meanExcess_div}
  \leanok
  If $F$ is antitone and regularly varying with index $-\alpha$, $\alpha > 1$, then
  $e(u)/u \to 1/(\alpha - 1)$ as $u \to \infty$. This matches the mean
  $\beta(u)/(1 - \xi) = u/(\alpha - 1)$ of the limiting GPD with $\xi = 1/\alpha$ and
  $\beta(u) = u/\alpha$ [EKM, \S3.4], [MFE, \S5.2].
\end{theorem}
\begin{proof}
  \leanok
  $e(u)/u = \big( \int_u^\infty F(t) \, \ud t \big) / \big( u \, F(u) \big)$, so this is exactly
  Karamata's theorem (Theorem~\ref{thm:karamata}).
\end{proof}
